\documentclass[11pt]{article}

\usepackage[a4paper, margin=1in]{geometry}
\usepackage{amsmath}
\usepackage{booktabs}
\usepackage{graphicx}
\usepackage{hyperref}
\usepackage{microtype}
\usepackage{natbib}
\usepackage{authblk}

\hypersetup{
    colorlinks=true,
    linkcolor=blue,
    citecolor=blue,
    urlcolor=blue,
}

\title{\textbf{Peak vs.\ Non-Peak Performance of Commercial LLM APIs:\\
A 28-Day Measurement Study}}

\author[1]{Jaekyung Noh}
\affil[1]{Independent Researcher \\ \texttt{silverbreak1@gmail.com}}

\date{May 2026}

\begin{document}

\maketitle

\begin{abstract}
Commercial large language model (LLM) APIs are shared infrastructure, yet little is known
about how server load during peak business hours affects latency, throughput, and answer
quality in practice.
We benchmark four providers---OpenAI (\texttt{gpt-4o-mini}), Anthropic
(\texttt{claude-haiku-4-5}), Google (\texttt{gemini-flash-lite}), and Alibaba Qwen
(\texttt{qwen3.5-flash})---twice daily for 28 days using 30 MMLU questions, measuring
time-to-first-token (TTFT), tokens-per-second (TPS), and answer accuracy.
Peak sessions run at 10:00\,ET (within the 8\,AM--2\,PM ET peak band \citep{anthropic2026promotion});
non-peak sessions run at 22:00\,ET.
We find that TTFT increases significantly during peak hours for US-hosted providers
(+4\% to +37\%), while Qwen---served from an Alibaba Cloud US endpoint---is
counter-intuitively \emph{faster} during US peak hours ($-7\%$ TTFT, $+9\%$ TPS),
consistent with the hypothesis that US peak corresponds to Asia off-peak.
Answer accuracy shows no statistically significant change for any provider across
conditions, suggesting that output quality is decoupled from infrastructure load at
the scale tested.
\end{abstract}

\section{Introduction}

Large language model APIs have become commodity infrastructure for production applications,
yet their service-level properties under varying load conditions remain largely undocumented
from the consumer's perspective.
Unlike traditional web services whose latency characteristics are well-studied, LLM inference
involves substantial per-request compute that amplifies the effect of concurrent demand.

Traffic on commercial AI APIs follows predictable diurnal patterns driven by North American
business hours.
If providers do not over-provision sufficiently, practitioners who call APIs during peak hours
may receive degraded latency---increasing costs for time-sensitive applications such as
real-time agents or interactive chat.
Whether throughput degradation also affects output \emph{quality} (answer accuracy) is an
open question with direct implications for reliability-critical deployments.

This paper reports the results of a 28-day empirical study comparing peak-hour and
off-peak-hour performance across four major LLM API providers on a standardized MMLU
benchmark suite.
Our primary contributions are:
\begin{enumerate}
    \item A reproducible measurement framework for longitudinal LLM API benchmarking.
    \item Provider-specific quantification of peak-vs-nonpeak latency and throughput effects.
    \item Evidence that MMLU accuracy is robust to infrastructure load conditions.
    \item An anomalous finding for a non-US-headquartered provider suggesting
          inter-regional load balancing.
\end{enumerate}

\section{Methodology}

\subsection{Providers and Models}

We evaluated four providers and their smallest fast-inference models available at the time
of the study (Table~\ref{tab:models}).
All models were prompted with a fixed system prompt instructing single-letter responses
(\texttt{A}--\texttt{D}), with \texttt{temperature=0} for reproducibility.
Streaming inference was used for all calls to measure TTFT directly.

\begin{table}[h]
\centering
\caption{Models under evaluation.}
\label{tab:models}
\begin{tabular}{lll}
\toprule
\textbf{Provider} & \textbf{Model} & \textbf{Endpoint region} \\
\midrule
OpenAI    & \texttt{gpt-4o-mini-2024-07-18}        & US East \\
Anthropic & \texttt{claude-haiku-4-5-20251001}     & US East \\
Google    & \texttt{gemini-3.1-flash-lite-preview}  & US (GCP) \\
Qwen      & \texttt{qwen3.5-flash}                  & US (Alibaba Cloud) \\
\bottomrule
\end{tabular}
\end{table}

\subsection{Question Set}

We selected 30 questions from the Massive Multitask Language Understanding (MMLU) benchmark
\citep{hendrycks2021measuring}, sampling across subjects and difficulty levels.
Each question was repeated three times per session to reduce variance, yielding 90 trials
per provider per session (30 questions $\times$ 3 repeats).
Accuracy was scored by extracting the first valid option letter from the model response
using a cascade of regular expressions targeting XML-tagged answers, explicit
``Answer:'' labels, and trailing isolated letters.

\subsection{Experimental Schedule}

Experiments ran from 2026-03-30 to 2026-04-26 (28 days).
Two sessions were scheduled daily in the America/New\_York timezone:
\begin{itemize}
    \item \textbf{Peak}: 10:00\,ET --- within the 8\,AM--2\,PM ET peak band \citep{anthropic2026promotion}.
    \item \textbf{Non-peak}: 22:00\,ET --- US East Coast late night.
\end{itemize}
Sessions were managed by a Python scheduler (APScheduler) running as a launchd daemon on
a Mac Mini with a residential internet connection.
Progress was checkpointed per task to support resumption after interruptions.

\subsection{Metrics}

\begin{itemize}
    \item \textbf{TTFT} (time-to-first-token, seconds): wall-clock time from request
          dispatch to receipt of the first streamed token.
    \item \textbf{TPS} (tokens per second): output token count divided by generation time
          after the first token.
    \item \textbf{Accuracy}: proportion of trials where the parsed answer matches the
          correct MMLU key.
\end{itemize}

\subsection{Statistical Analysis}

All pairwise comparisons use the Mann--Whitney U test (two-sided), which makes no
distributional assumptions and is robust to the heavy-tailed latency distributions observed.
Effect sizes are reported as median differences and percent change relative to the non-peak
median (mean for accuracy). Significance threshold: $\alpha = 0.05$.

\section{Results}

\subsection{Data Completeness}

Of 56 scheduled slots (28 days $\times$ 2 periods), 54 were executed; 3 non-peak slots
were missed due to daemon downtime.
Most executed slots were partial: API errors (primarily Google returning HTTP~429/503)
caused premature termination in 45 of 54 slots.
A \emph{complete} slot is defined as exactly 360 rows
(30 questions $\times$ 3 repeats $\times$ 4 providers).
Only 7 slots were complete (3 non-peak, 4 peak).
All primary results below use these 7 complete slots (1,890 rows total);
see Section~\ref{sec:limitations} for discussion of this constraint.

\subsection{Latency: Time-to-First-Token}

\begin{table}[h]
\centering
\caption{TTFT comparison across conditions (complete slots only; $n_\text{peak}=360$,
$n_\text{nonpeak}=270$ per provider; medians reported).
Significance: $^{***}p<0.001$, $^{*}p<0.05$, \textit{ns}=not significant.}
\label{tab:ttft}
\begin{tabular}{lrrrrr}
\toprule
\textbf{Provider} & \textbf{Non-peak (s)} & \textbf{Peak (s)}
  & \textbf{$\Delta$ (s)} & \textbf{$\Delta$ (\%)} & \textbf{Sig.} \\
\midrule
OpenAI    & 0.873 & 0.910 & $+0.037$ & $+4.3\%$  & $^{*}$   \\
Anthropic & 0.835 & 0.971 & $+0.136$ & $+16.3\%$ & $^{***}$ \\
Google    & 1.721 & 2.355 & $+0.634$ & $+36.8\%$ & $^{***}$ \\
Qwen      & 3.698 & 3.445 & $-0.253$ & $-6.8\%$  & $^{***}$ \\
\bottomrule
\end{tabular}
\end{table}

Three of the four providers exhibit significantly higher TTFT during peak hours
(Table~\ref{tab:ttft}).
Google shows the largest absolute and relative increase (+37\%), while OpenAI shows the
smallest (+4\%).
Qwen is the sole outlier: TTFT is \emph{lower} during peak hours, a pattern we discuss
in Section~\ref{sec:discussion}.

\subsection{Throughput: Tokens per Second}

\begin{table}[h]
\centering
\caption{TPS comparison (complete slots only; medians).}
\label{tab:tps}
\begin{tabular}{lrrrrr}
\toprule
\textbf{Provider} & \textbf{Non-peak} & \textbf{Peak}
  & \textbf{$\Delta$} & \textbf{$\Delta$ (\%)} & \textbf{Sig.} \\
\midrule
OpenAI    & 54.9  & 53.6  & $-1.4$  & $-2.5\%$ & \textit{ns} \\
Anthropic & 139.5 & 139.1 & $-0.5$  & $-0.3\%$ & \textit{ns} \\
Google    & 221.6 & 226.9 & $+5.3$  & $+2.4\%$ & \textit{ns} \\
Qwen      & 171.3 & 187.3 & $+15.9$ & $+9.3\%$ & $^{***}$    \\
\bottomrule
\end{tabular}
\end{table}

TPS differences are not statistically significant for the three US-hosted providers
(Table~\ref{tab:tps}).
Qwen again shows the opposite pattern: generation throughput is significantly higher during
US peak hours.

\subsection{Answer Accuracy}

\begin{table}[h]
\centering
\caption{MMLU accuracy (mean correct, complete slots only).}
\label{tab:accuracy}
\begin{tabular}{lrrrl}
\toprule
\textbf{Provider} & \textbf{Non-peak} & \textbf{Peak}
  & \textbf{$\Delta$} & \textbf{Sig.} \\
\midrule
OpenAI    & 0.493 & 0.475 & $-0.018$ & \textit{ns} \\
Anthropic & 0.833 & 0.817 & $-0.017$ & \textit{ns} \\
Google    & 0.867 & 0.864 & $-0.003$ & \textit{ns} \\
Qwen      & 0.822 & 0.808 & $-0.014$ & \textit{ns} \\
\bottomrule
\end{tabular}
\end{table}

No provider shows a statistically significant change in accuracy between peak and non-peak
conditions (Table~\ref{tab:accuracy}, all $p > 0.5$).
The sub-50\% accuracy of OpenAI's \texttt{gpt-4o-mini} reflects this model's sensitivity to
strict single-letter output formatting; format-failure rates range from 1.5\% to 2.8\%.

\section{Discussion}
\label{sec:discussion}

\paragraph{US-hosted providers: TTFT scales with load.}
OpenAI, Anthropic, and Google all exhibit elevated TTFT during peak hours, consistent with
increased queuing or scheduling overhead under higher concurrent demand.
The effect ranges from modest for OpenAI ($\approx$4\%) to substantial for Google
($\approx$37\%), suggesting different capacity margins across providers.
Critically, this latency increase does not propagate to generation speed or accuracy,
implying that once the model begins generating, execution is largely isolated from
request-level congestion.

\paragraph{Qwen anomaly: cross-regional load balancing.}
Qwen's US endpoint (\texttt{dashscope-us.aliyuncs.com}) runs on Alibaba Cloud
infrastructure.
US business peak hours (10:00\,ET) correspond to approximately 23:00--01:00 the following
day in China Standard Time---well outside Chinese business hours.
The observed lower TTFT and higher TPS during US peak hours is consistent with the hypothesis
that Alibaba's global scheduler shifts load to US-region GPUs when Asian demand is low,
effectively benefiting US users during their own peak.
This finding suggests that infrastructure geography and inter-regional load balancing may
produce effects opposite to naive expectations for non-US-headquartered providers.

\paragraph{Accuracy robustness.}
The absence of accuracy degradation during peak hours is practically important.
Applications relying on LLM APIs for classification or reasoning can expect consistent
output quality regardless of the time of day.
Operators optimizing purely for quality need not schedule around traffic patterns.
Those optimizing for latency-sensitive workloads (real-time agents, interactive chat) should
budget for TTFT increases of 4--37\% during US peak hours.

\section{Limitations}
\label{sec:limitations}

\paragraph{Limited complete-slot coverage and Google data gaps.}
Of 54 executed slots, only 7 met the completeness threshold of 360 rows.
The primary cause was Google Gemini API rate-limiting (HTTP~429/503 errors), which
terminated experiments mid-session, resulting in entire Google response rows being absent
rather than malformed.
This disproportionately affected non-peak slots and may introduce selection bias.
Analyses are restricted to complete slots to minimize this bias, but the sample covers only
a subset of the 28-day window and reduces statistical power.

\paragraph{Single time point per condition.}
Each daily session is a single snapshot (10:00 or 22:00\,ET).
Intra-period variability (e.g., 9:00 vs.\ 13:00 during peak) is not captured.

\paragraph{Residential network variability.}
All measurements were made from a Mac Mini on a residential internet connection.
Network jitter contributes to TTFT variance and cannot be fully separated from server-side
latency effects.

\paragraph{Small question set and statistical power.}
30 MMLU questions are insufficient to detect small accuracy effects ($<$2\%) or
subject-level differences. A larger question set would improve sensitivity.

\paragraph{Non-comparable generation speed across providers.}
Qwen (\texttt{qwen3.5-flash}) returns extended chain-of-thought reasoning by default,
producing responses substantially longer than those of other providers for the same prompt.
As a result, TPS figures for Qwen reflect a different output regime and direct
cross-provider speed comparisons are not meaningful.
Within-provider peak vs.\ non-peak TPS comparisons remain valid.

\paragraph{Model versioning.}
Providers may silently update model weights or infrastructure during the study period.
Observed trends reflect the combined effect of load patterns and any within-study
infrastructure changes, which we did not control for.

\section{Conclusion}

We conducted a 28-day longitudinal benchmark of four commercial LLM APIs at US business-hour
peak and off-peak times.
US-hosted providers (OpenAI, Anthropic, Google) show significantly elevated TTFT during peak
hours (4--37\%), with no corresponding degradation in generation throughput or answer
accuracy.
Qwen, served from a US-region Alibaba Cloud endpoint, shows the opposite pattern---faster
TTFT and higher TPS during US peak---consistent with cross-regional load balancing.

These findings suggest that practitioners should expect latency increases for latency-critical
applications during US business hours but can rely on consistent output quality regardless of
time of day.
Future work should extend the study period, increase question coverage, and instrument
server-side metrics to corroborate the load-balancing hypothesis.

\section*{Data and Code Availability}

The benchmark framework, collected CSVs, and analysis scripts are available at
\url{https://github.com/silverbreak1/llm_trafic_router}.

\bibliographystyle{plainnat}
\bibliography{references}

\end{document}
